% !TEX root = chapter-standalone.tex

\section{Functors}
\label{sec:functors}
\label{sec:functors-intro}
\linkvideo{spring2021-functors:bigpic} % Big picture, zipping
\linkvideo{spring2021-functors:semi-and-fun:fun-def} % Definition of functor




Any given category \CatC is a ``compositional world'': inside \CatC we have objects, morphisms between them, and a way to talk about composing morphisms.
Now we will zoom out a level, and consider different categories -- different worlds -- simultaneously, and how to relate them to each other.

The most basic notion of how to relate one category to another is given by the concept of a \emph{functor}.

Just like a morphism
%
\begin{equation}
    \mora \colon \Obja \mto \Objb
\end{equation}
%
relates a source object to a target object inside a given category, a functor
%
\begin{equation}
    \funa \colon \CatC \fto \CatD
\end{equation}
%
relates a source category to a target category. In fact, we will see in the next chapters that we can view functors as morphisms inside a \emph{category of categories}.

% \todotextjira{118}{@Andrea: add introduction here, with several examples of translations}

Here is the formal definition of a functor; we'll follow it with many examples to illustrate. 


\begin{ctdefinition}[Functor]
    \label{def:functor}
    Given \SY{categories}~\CatC and~\CatD, a \maindef{functor}~$\funa\colon \CatC\fto \CatD$ from~\CatC to~\CatD is:

    \constit

    \begin{enumerate}
        \item A function
              \begin{equation}
                  \funaob \colon \ObC \to \ObD.
              \end{equation}
        \item For every pair of objects~$\Obja, \Objb\setin \Obof\CatC$, a function
              \begin{equation}
                  \funamor\colon \HomSet{\CatC}{\Obja}{\Objb} \to \HomSet{\CatD}{\funaob(\Obja)}{\funaob(\Objb)}.
              \end{equation}
    \end{enumerate}

    \condit

    \begin{enumerate}
        \item The function $\funamor$ is compatible with the composition operations in the source and target category:
              \begin{equation}
                  \prfperiod{
                      \mora\colon \Obja \mtoin{\CatC} \Objb
                  }{\quad}{
                      \morb\colon \Objb \mtoin{\CatC} \Objc
                  }{
                      \funamor(\mora \mthenof{\CatC} \morb)=\funamor(\mora) \mthenof{\CatD} \funamor(\morb)
                  }
              \end{equation}
       \item The function $\funamor$ is compatible with identity morphisms in the source and target category:
       \begin{equation}
        \label{eq:functor-condition}
        \funamor(\catidat\Obja)=\catidat{\funaob(\Obja)}
    \end{equation}
    for all objects $\Obja$ in $\CatC$.

    \end{enumerate}
\end{ctdefinition}


\begin{example}[Powerset functor]
    \label{ex:powerset_functor}
    We define a \SY{functor}~$\funa\colon \Set \fto \Set$ which maps each set to its power set.
    In other words, $\funaob(\setA) = \powerset \setA$ for any set $\setA$.
    On the level of morphisms, given any function $\mora \colon \setA \mto \setB$, we define
    \begin{equation}\label{eq:powerset-function-def}
        \begin{aligned}
            \funamor(\mora)\colon \powersetof \setA & \to \powersetof \setB \\
            \setC                                   & \mapsto \makeset{\mora(\setCel) \mid \setCel \setin \setC}.
        \end{aligned}
    \end{equation}
    Here is a concrete illustration.
    Consider the two sets~$\setA = \makeset{\sfondue, \sbretzel, \schoco}$ and~$\setB = \makeset{\scheese,\sgrapes,\sapple}$.
    Applying the \SY{functor} to~\setA gives
    \begin{equation}\label{eq:powerset-functor-example-1a}
        \funaob(\setA)
        =
        \makeset{
            \Emptyset,
            \makeset{\sfondue},
            \makeset{\sbretzel},
            \makeset{\schoco},
            \makeset{\sfondue,\sbretzel},
            \makeset{\sfondue,\schoco},
            \makeset{\sbretzel,\schoco},
            \makeset{\sfondue,\sbretzel,\schoco}
        }
    \end{equation}
    and applying it to $\setB$ gives
    \begin{equation}\label{eq:powerset-functor-example-1b}
        \funaob(\setB)
        =
        \makeset{
            \Emptyset,
            \makeset{\scheese},
            \makeset{\sgrapes},
            \makeset{\sapple},
            \makeset{\scheese,\sgrapes},
            \makeset{\scheese,\sapple},
            \makeset{\sgrapes,\sapple},
            \makeset{\scheese,\sgrapes,\sapple}
        }.
    \end{equation}
    Furthermore, consider the map
    %
    \begin{equation}\label{eq:powerset-functor-example-2}
        \begin{aligned}
            \mora \colon \setA & \sto \setB, \\
            \sfondue           & \mapsto \scheese, \\
            \sbretzel          & \mapsto \sgrapes, \\
            \schoco            & \mapsto \sapple.
        \end{aligned}
    \end{equation}
    %
    This would for instance give~$\funamor(\mora)(\makeset{\sfondue,\sbretzel})=\makeset{\mora(\sfondue),\mora(\sbretzel)}=\makeset{\scheese,\sgrapes}$.

    Now let us check that $\funa$, so-defined, really is a functor.

    First let us check that it is compatible with composition.
    Consider functions~$\mora\colon \setA\fto \setB$,~$\morb\colon \setB \fto \setC$.
    On the one hand we have
    \begin{equation}\label{eq:powerset-functor-example-3}
        \begin{aligned}
            \funa(\morab)(\setC) & =\makeset{ \morb(\mora(\setCel))\mid \setCel \setin \setC},
        \end{aligned}
    \end{equation}
    and on the other hand
    \begin{equation}\label{eq:powerset-functor-example-4}
        \begin{aligned}
            (\funa(\mora)\mthen \funa(\morb))(\setC)
             & =\funa(\morb) (\makeset{ \mora(\setCel) \mid \setCel \setin \setC}) \\
             & =\makeset{g(\setDel)\mid \setDel \setin \makeset{\mora(\setCel)\mid \setCel\setin \setC}} \\
             & =\makeset{ \morb(\mora(\setCel))\mid \setCel \setin \setC}.
        \end{aligned}
    \end{equation}
    Second, let us check that $\funa$ is compatible with identity morphisms.
    We have
    \begin{equation}\label{eq:powerset-functor-example-5}
        \begin{aligned}
            \funa(\catidat\setA)(\setC) & = \makeset{\catidat\setA(\setCel)\mid \setCel \setin \setC} \\
                                        & =\catidat{\funa(\setC)}.
        \end{aligned}
    \end{equation}
\end{example}

\begin{remark}
\cref{ex:powerset_functor}, the powerset functor, is an example of an \emph{endofunctor}: a functor from a category to itself. 
\end{remark}

\begin{example}
    There is a \SY{functor}
    \begin{equation}\label{eq:list-functor-intro}
        \listfun\colon \Set \fto \Mon
    \end{equation}
    from the \SY{category of sets} to the \SY{category of monoids}, defined as follows.

    Given a set~\setA, the \SY{functor} returns a specific \SY{monoid}
    \begin{equation}\label{eq:list-functor-on-objects}
        \listfun(\setA)\definedas\tupp{\listsof{\setA},\emptylist_\setA, \listconcat}.
    \end{equation}

    Given a map~$\mapa\colon \setA \sto \setB$, we have
    % \begin{equation}
    %     \defmapcomma{
    %         \listfun(\mapa)
    %     }{
    %         \listsof{\setA}
    %     }{
    %         \to
    %     }{
    %         \listsof{\setB}
    %     }{
    %         \makelist{\setAel_1, \ldots, \setAel_n}
    %     }{
    %         \makelist{\mapa(\setAeln{1}), \ldots, \mapa(\setAeln{n})}
    %     }
    % \end{equation}
    \begin{equation}\label{eq:list-functor-on-morphisms}
        \defmapcomma{
            \listfun(\mapa)
        }{
            \listsof{\setA}
        }{
            \to
        }{
            \listsof{\setB}
        }{
            \makegenlist{ \setAel_i }
        }{
            \makegenlist{ \mapa(\setAeln{i}) }
        }
    \end{equation}
    which applies~$\mapa$ entry-wise in the list.
    The empty list in~$\listsof{\setA}$ is mapped to the empty list in~$\listsof{\setB}$.
        {}
    $\listfun$ is a \SY{functor}, because identity functions in $\Set$ are mapped to the corresponding identity morphisms on lists, and
    % \begin{equation}
    %     \begin{aligned}
    %         \listfun(\mapa \mthen \mapb)(\makelist{\setAeln{1},\ldots,\setAeln{n}}) & =\makelist{(\mapa \mthen \mapb)(\setAeln{1}), \ldots, (\mapa \mthen \mapb)(\setAeln{n})} \\
    %                                                                                 & =\listfun(\mapb)(\makelist{\mapa(\setAeln{1}), \ldots, \mapa(\setAeln{n})}) \\
    %                                                                                 & =(\listfun(\mapa)\mthen \listfun(\mapb))(\makelist{\setAeln{1}, \ldots, \setAeln{n}}).
    %     \end{aligned}
    % \end{equation}
    \begin{equation}\label{eq:list-functor-compositionality}
        \begin{aligned}
            \listfun(\mapa \mthen \mapb)(\makegenlist{\setAeln{i}}) & =\makegenlist{ (\mapa \mthen \mapb)(\setAeln{i})} \\
                                                                    & =\makegenlist{  \mapb ( \mapa (\setAeln{i}))} \\
                                                                    & =\listfun(\mapb)(\makegenlist{\mapa(\setAeln{i})}) \\
                                                                    & =(\listfun(\mapa)\mthen \listfun(\mapb))(\makegenlist{\setAeln{i}}).
        \end{aligned}
    \end{equation}
\end{example}

% To see that this is indeed a \SY{functor}, notice that~$\funidC(\morab)$ is equal to~$\funidC(\mora)\mthen \funidC(\morb)$ because they are both equal to~$\morab$.
% Moreover, identities are preserved because~$\funidC(\catidof\CatC)=\catidof\CatC$.



\subsection{Functors generalize monotone functions.}
\label{sec:posetsarecats}

\linkvideo{spring2021-tradeoffs:tradeoffs:orders:preorder-as-cat} % Pre-order as a category
\linkvideo{spring2021-tradeoffs:tradeoffs:orders:preorder-poset} % The skeleton of a pre-order is a poset

Recall that a single \SY{poset}~$\posAdefinition$ can be viewed as a category -- let's call it $\poscat{\posA}$ -- in which each element of the poset is an object, and there is exactly one morphism ~$\posela \mto \poselb$ if and only if~$\posela \posAleq \poselb$ (\cref{def:poscat}), and otherwise there is no such morphism.

% \subsection{Monotone maps are functors}

\linkvideo{spring2021-functors:semi-and-fun:mon-functions:mon-fun-as-func} % Monotone functions as functors

\begin{lemma}
    \label{lem:posetfunctor}
    Given posets $\posA$ and  $\posB$, a \SY{monotone function} $\posA \mto \posB$ is the ``same thing'' as a \SY{functor} $\poscat{\posA} \mto \poscat{\posB}$.
\end{lemma}


%\begin{proof}
%    We start by specifying the \SY{functor}~$\funa$ and two posetal categories~$\poscat{\posA}$ and~$\poscat{\posB}$.
%    We first specify the action of~$\funa$ on objects (elements of a \SY{poset} considered as objects of the posetal category) and on morphisms (order relations considered as morphisms of the posetal category).
%    A \SY{monotone function} maps each element of a poset~$\posela\setin \posAset$ to~$\funa(\posela) \setin \posBset$, and it guarantees that for every~$\posela,\poselb\setin \posAset$, if~$\posela\posAleq \poselb$ then~$\funa(\posela)\posBleq \funa(\poselb)$.
%    We now need to check the two conditions that a \SY{functor} must satisfy.
%    First, consider the \SY{identity morphism} for~$\posela\setin \posAset$, namely~$\posela \posAleq \posela$.
%    The application of the map~$\funa$ results in the condition~$\funa(\posela)\posBleq \funa(\posela)$, which is the \SY{identity morphism} on~\posB.
%    Second, morphisms~$\posela\posAleq \poselb$ and~$\poselb\posAleq \poselc$ in~\posA, by applying the map~$\funa$ to the morphism composition~$\posela\posAleq \poselc$ we obtain~$\funa(\posela)\posBleq \funa(\poselc)$, which is the same as the composition of~$\funa(\posela)\posBleq \funa(\poselb)$ and~$\funa(\poselb)\posBleq \funa(\poselc)$.
%\end{proof}

\todomistake{J: I find the previous version of the lemma statement (which I've rewritten) and also the proof (which I have commented out for now), confusing.
    In particular, I would not use the word ``isomorphic'' in the statement of the lemma, since I think this cannot be made into a precise/correct statement.
}

%\devel{\includepdf[scale=0.8,pages={25-26},nup=1x3,frame,pagecommand={}]{ACT4E-06-posets.pdf}} % preorder as category


\subsection{Functors generalize monoid morphisms.}
\label{sec:functors-generalize-monoid-morphisms}

\linkvideo{spring2021-functors:semi-and-fun:ex-semigroup-semifun} % Semigroup morphisms as semi-functors


Recall as well that given a monoid $\monoidA$, we can view it as a category -- let's call it $\poscat{\monoidA}$ -- where the elements of the monoid play the role of morphisms.

From this point of view, a functor corresponds to a homomorphism of monoids.

\begin{lemma}
    \label{lem:monoid-hom-as-functor}
    Given monoids $\monoidA$ and $\monoidB$, a \SY{monotone function} $\monoidA \mto \monoidB$ is the ``same thing'' as a \SY{functor} $\poscat{\monoidA} \mto \poscat{\monoidB}$.
\end{lemma}


\subsection{Functors generalize graph homomorphisms.}


We have seen that given a (directed multi-) graph $\graphA=\tupp{\vertices,\arcs, \source,\target}$, we can turn it into a category -- let's call the result $\poscat{\graphA}$ -- where the vertices of the graph become the objects of $\poscat{\graphA}$ and the morphism in $\poscat{\graphA}$ are paths in $\graphA$.

\begin{lemma}
    \label{lem:graph-hom-to-functor}
    Given graphs $\graphA$ and $\graphB$, a \SY{graph homomorphism} $\graphA \mto \graphB$ induces a \SY{functor} $\poscat{\graphA} \mto \poscat{\graphA}$.
\end{lemma}

\todotext{J: check if there are however functors $\poscat{\graphA} \mto \poscat{\graphA}$ which are *not* induced by graph homomorphisms... my gut feeling is `yes'. And I think this is an interesting point to come back to when discussing adjunctions}


\todojira{620}{\alphubel: @JL: let's discuss the ``image'' of a \SY{functor} somewhere}
